\section{External premises and source provenance}\label{sec:provenance}
The remaining non-elementary external theorem premises are exactly
Coladangelo v1 Theorem 7.1, the finite-projective tensor clause of
Theorem 7.2, Proposition 10.9, and the standard finite-dimensional
real-analytic stable/unstable manifold theorem. The first three are stated
in External premise~\ref{prem:Coladangelo} and refer to one fixed v1 PDF,
whose identity is recorded in Appendix~\ref{app:audits}. They share proof
ancestry: Theorem 7.2 uses the scalar-certificate theory, while Proposition
10.9 also uses the source's finite-approximation and equality arguments.
Multiple AI-generated branches that consume these premises are not
independent confirmations of them.

The parameterization-method article~\cite{CabreFontichLlave2003} supplies
a standard reference for analytic regularity, with the stable/unstable
specialization described in its Theorem 1.1 discussion. The textbook
\cite{Shub1987} is a reference for the classical local theory. The frozen
provenance audit inspected the analytic article; it verified the textbook's
publisher metadata but did not successfully obtain and inspect its full
text. This is not represented as a new independent proof of the general
theorem. Its concrete applicability here---analyticity, invertibility,
hyperbolicity and direction counts---is established in
Lemma~\ref{lem:tails}.

Pauwels's preprint and its associated formal development are comparison
sources, not premises of the final theorem. In particular, their formal
statements are not a Lean certification of the residual zero or of this
high-precision enclosure. P\'al--V\'ertesi is historical context and
construction ancestry. No reported floating-point optimum or conjectured
optimality from that work is assumed.

Douglas's v1.0.0 Zenodo archive~\cite{Douglas2026Zenodo} is cited as a
historical source for the candidate constant and matching picture, and the
later arXiv paper~\cite{Douglas2026Arxiv} for related subsequent work. These
sources are included for historical and methodological provenance and are not
used as premises of the theorem proved here.

The R1 candidates are the two specific rational objects released by
Mghirbi in v2.0.0~\cite{Mghirbi2026data}. Released-byte identity and exact
validity have been independently certified. Complete numerical-generation
ancestry before that release has not been reconstructed. This distinction
matters: the proof of Proposition~\ref{prop:R1} accepts the explicit
projectors and polynomial identity after exact checking, and does not need
the original solver trajectory to be correct or recoverable. The lower
decimal printed in the frozen companion preprint is inaccurate in its last
reported digits; the rational witness and its recomputed interval in this
paper are the authority. Later source versions do not replace the frozen
upper certificate.

The source audit records other local textual errors, including a false
sentence asserting that shifting an irreducible Jacobi matrix by a scalar
identity makes all its entries positive. Such a shift cannot change the
zero corner entries of a tridiagonal matrix. The actual weaker fact used
here is positivity on a connected component, established by absolute-value
maximization and zero propagation in Appendix~\ref{app:compactness}; no
false source sentence is imported. Appendix~\ref{app:audits} records this
and the remaining citation caveats.

The Douglas sources are cited for historical and methodological provenance
rather than as mathematical premises of the final proof chain. The
repository-level history records that the earlier matching picture informed
later question formulation. This distinction between historical influence
and proof dependence is maintained throughout the final assembly. The
assembly used no project-specific material outside the supplied archive;
fresh web access was limited to bibliographic metadata checks, and no proof
step was imported from it.
